\documentclass{beamer}
\usefonttheme{professionalfonts}
%\usetheme{Rochester}
%\usetheme{AnnArbor}
\usetheme{Szeged}
\useoutertheme{infolines}
\usecolortheme{dolphin}
\usepackage{graphicx}
% \usepackage[english]{babel}
\usepackage{hyperref}
\usepackage{multicol}
\usepackage{color}
\usepackage{xcolor}
%\usepackage{smartdiagram}


%\definecolor{carmine}{rgb}{0.59, 0.0, 0.09}
\usepackage{amsmath} 

\usepackage{graphicx}
\usepackage{tikz}
\usetikzlibrary{shapes.callouts}
\usepackage{multirow}

\usepackage[utf8]{inputenc}

\usepackage{enumerate}
\usepackage[shortlabels]{enumitem}
\usepackage{amsmath,amsthm,amssymb,amsfonts,mathrsfs,latexsym,stmaryrd}
\usepackage{tcolorbox}

\usepackage{color} 
\usepackage{amssymb, amsmath, amsbsy}
\usepackage{adjustbox,lipsum}
\usepackage{placeins}

\usepackage{textpos}
\usepackage{wrapfig}
\usepackage[spanish, es-tabla]{babel}

\usepackage{caption}
\setlength{\parindent}{0pt}
\usepackage[hyperfootnotes=false]{hyperref}

\usepackage{url}
\usepackage{float}
\usepackage{tikz}
\usetikzlibrary{patterns}


\title[ICS-162]{Econometría ICS-294\\ \small{Clase 8: Propiedades estadísticas del estimador MCO}}

\author{Jocelyn Tapia Stefanoni}
\institute[USM]
{Universidad Técnica Federico Santa María\\ % Your institution for the title page
	\medskip
	\textit{} % Your email address
	
}
\logo{

\includegraphics[width=3cm]{figuras/logoUSM.png}
}
\date{}
\newcommand\Fontvi{\fontsize{9}{7.2}\selectfont}
\begin{document}
	
\begin{frame}
	\titlepage % Print the title page as the first slide


\end{frame}
%%%%%%%%%%%%%%%%%%%%%%%%%%%%%%%%%%%%%%%%%%%%%%%%%%%%%%%%%%%%%%%%%%%%%%%%%%%
\logo{}






\begin{frame}
\frametitle{Estimador MCO de $\sigma^2$}
La matriz de varianzas-covarianzas viene dada por:
\[
\text{Var}(\hat{\beta}|X) = \sigma^2(X^{\top}X)^{-1}
\]
Notar que $(X^{\top}X)^{-1}$ es una matriz que se observa en los datos.

Pero, $\sigma^2$ es la varianza del error poblacional. NO es observable.

\end{frame}

\begin{frame}
\frametitle{Estimador MCO de $\sigma^2$}



El estimador $\hat{\sigma}^2$ de $\sigma^2$:

 $$\hat{\sigma}^2 = \frac{\sum_{i=1}^{n} \hat{u}_i^2 }{n-k-1}= \frac{SSR}{n - k - 1}$$


\begin{itemize}
    \item  El término n - k - 1  son los grados de libertad (gl) para el problema general de
MCO con n observaciones y k variables independientes.
\item  Como en un modelo de regresión con k variables independientes y un intercepto hay k + 1 parámetros.
\item gl =  (observaciones) - (cantidad de parámetros estimados).
\item  Bajo los supuestos del teorema de Gauss-Markov se cumple que $E[\hat{\sigma}^2|X] = \sigma^2$.
\end{itemize}

\end{frame}

% \begin{frame}{}
    
% \end{frame}
% Éstas pueden escribirse como Σ i1 n
% u ˆ i  0 y
% Σ i1 n
% xij u ˆ i  0, donde j  1, 2, …, k. Por tanto, en la obtención de las estimaciones de MCO,
% a los residuales de MCO se les impone k  1 restricciones. Esto significa que, dados n  (k 
% 1) residuales, los restantes k  1 residuales se conocen: sólo hay n  (k  1) grados de libertad
% en los residuales. (Esto se puede comparar con los errores ui, que tienen n grados de libertad en
% la muestra.)
 \begin{frame}
\frametitle{Varianza Estimada de Cada Estimador $\hat{\beta}_j$}

\[
\hat{\text{Var}}(\hat{\beta}|X) = \hat{\sigma}^2(X^{\top}X)^{-1}
\]
Matriz de $(k + 1) \times (k + 1)$ de varianzas-covarianzas, tal que:
\begin{itemize}
    \item El elemento de la diagonal $a_{1,1}$ es $\hat{\text{Var}}(\hat{\beta}_0)$.
    \item El elemento de la diagonal $a_{2,2}$ es $\hat{\text{Var}}(\hat{\beta}_1)$.
    \item Generalizando: el elemento de la diagonal $a_{j+1,j+1}$ es $\hat{\text{Var}}(\hat{\beta}_j)$.
\end{itemize}
\end{frame}
%%%%%%%%%%%%%%%%%%%%%%%%%%%%%%%%%%%%%%%%%%
% \begin{frame}
% \frametitle{Ejemplo}
% Estimamos el siguiente modelo con dos variables:
% \[
% \text{salario}_i = \beta_0 + \beta_1 \text{educ}_i + \beta_2 \text{IQ}_i + u_i
% \]
% \end{frame}
%%%%%%%%%%%%%%%%%%%%%%%%%%%%%%%%%%%%%%%%%%%
% \begin{frame}
% \frametitle{Ejemplo}
% Estimamos el siguiente modelo con dos variables:
% \[
% \text{salario}_i = \beta_0 + \beta_1 \text{educ}_i + \beta_2 \text{IQ}_i + u_i
% \]

% \begin{figure}
%     \centering
%     \includegraphics[width=0.8\linewidth]{clase4/images/ej.png}
% \end{figure}

% \end{frame}
%%%%%%%%%%%%%%%%%%%%%%%%%%%%%%%%%%%%%%%%%%%
\begin{frame}
\frametitle{Varianza del Estimador MCO: Ejemplo}


\begin{itemize}
    \item Supuestos clásicos $\Rightarrow \hat{V}(\hat{\beta}) = \hat{\sigma}^2(X^{\top}X)^{-1}$.
    \item De los datos obtenemos:
    \[
    (X^{\top}X)^{-1} =
    \begin{bmatrix}
    0.05987 & -0.00177 & -0.00035 \\
    -0.00177 & 0.00030 & -0.00002 \\
    -0.00035 & -0.00002 & 0.00001 \\
    \end{bmatrix}; \]
    
\[\hat{\sigma}^2 = \frac{SSR}{n-k-1} =141925
    \]
%     n=28\\
    
% $SSR=\sum_{i=1}^{n} \hat{u}_i^2=3.548.125$
    
\end{itemize}
\end{frame}
% %%%%%%%%%%%%%%%%%%%%%%%%%%%%%%%%%%%%%%%%%%%
% \begin{frame}
% \frametitle{Varianza del Estimador MCO: Ejemplo}
% \begin{itemize}
%     \item Supuestos clásicos $\Rightarrow \hat{V}(\hat{\beta}) = \hat{\sigma}^2(X^{\top}X)^{-1}$.
%     \item De los datos obtenemos:
%     \[
%     (X^{\top}X)^{-1} =  \begin{bmatrix}
%     0.05987 & -0.00177 & -0.00035 \\
%     -0.00177 & 0.00030 & -0.00002 \\
%     -0.00035 & -0.00002 & 0.00001 \\
%     \end{bmatrix}
%     \]

%     $ \quad \hat{\sigma}^2 = 141925$
%     \item Entonces, podemos calcular:
%     \[
%     \hat{V}(\hat{\beta}) = \hat{\sigma}^2(X^{\top}X)^{-1} =
%     \begin{bmatrix}
%     8497.57 & -250.80 & -49.04 \\
%     -250.80 & 42.90 & -3.22 \\
%     -49.04 & -3.22 & 0.91 \\
%     \end{bmatrix}
%     \]
% \end{itemize}
% \end{frame}
 \begin{frame}
 \frametitle{Varianza del Estimador MCO: Ejemplo}
\begin{itemize}
    \item Supuestos clásicos $\Rightarrow \hat{V}(\hat{\beta}) = \hat{\sigma}^2(X^{\top}X)^{-1}$.
    \item De los datos obtenemos:
    \[
    (X^{\top}X)^{-1} =
    \begin{bmatrix}
    0.05987 & -0.00177 & -0.00035 \\
    -0.00177 & 0.00030 & -0.00002 \\
    -0.00035 & -0.00002 & 0.00001 \\
    \end{bmatrix}; \quad \hat{\sigma}^2 = 141925
    \]
    \item Entonces, podemos calcular:
    \[
    \hat{V}(\hat{\beta}) = \hat{\sigma}^2(X^{\top}X)^{-1} =
    \begin{bmatrix}
    8497.57 & -250.80 & -49.04 \\
    -250.80 & 42.90 & -3.22 \\
    -49.04 & -3.22 & 0.91 \\
    \end{bmatrix}
    \]
    \item En la diagonal principal, tienen la varianza de cada uno de los estimadores.
      \item Ejemplo: $\hat{V}(\hat{\beta}_1) = 42.6$. Es decir, los $\hat{se}(\hat{\beta}_1) = \sqrt{\hat{V}(\hat{\beta}_1)} = 6.53$
\end{itemize}
\end{frame}

\begin{frame}
\frametitle{El Caso Particular de Regresión con una Sola Variable Independiente}
En el modelo de regresión simple $y = \beta_0 + \beta_1 x + u$, sigue cumpliéndose lo anterior.
\[
\hat{Var}(\hat{\beta}_1) = {\frac{{\hat{\sigma}}^2}{\sum_{i=1}^{n}(x_i - \bar{x})^2}}
\]
\begin{itemize}
    \item Si $x$ tiene más variabilidad, hace que $\hat{\beta}$ sea más preciso (tengo más información).
    \item Si la varianza del residuo es más grande, mi estimador es menos preciso. Mi modelo no explica mucho de $y$.
\end{itemize}
 \end{frame}

% \begin{frame}[shrink=10]
% \frametitle{Varianza del Estimador MCO}
% \begin{itemize}
%     \item Intuición de la varianza del $\hat{\beta}_1$ (elemento $a_{22}$ de la diagonal):
%     \[
%     \hat{Var}(\hat{\beta}_1|X) = \frac{\hat{\sigma}^2}{SST_1(1 - R^2_1)}
%     \]
%     \item Donde:
%     \begin{itemize}
%         \item $SST_1 = \sum_{i=1}^{n}(x_{i1} - \bar{x}_1)^2 = n \cdot \text{Var}(x_1)$ es la suma total de los cuadrados de $x_1$.
%         \item $R^2_1$: el coeficiente de determinación resultante de la regresión de $x_1$ (como variable dependiente) en todas las otras variables independientes. Intuitivamente, indica cuánto de la variabilidad de $x_1$ se puede explicar por $x_2, x_3, \ldots$.
%     \end{itemize}
%     \item Ejemplo:
%     \begin{itemize}
%         \item Modelo: $y_i = \beta_0 + \beta_1x_{1i} + \beta_2x_{2i} + \beta_3x_{3i} + u_i$.
%         \item $R^2_1$ es el coeficiente de determinación de la regresión MCO auxiliar:
%         \[
%         x_{1i} = \alpha_0 + \alpha_2x_{2i} + \alpha_3x_{3i} + e_i
%         \]
%         $R^2_1$ indica cuánto de la variable $x_1$ es explicado por las otras variables independientes.
%     \end{itemize}
% \end{itemize}
% \end{frame}

% \begin{frame}[shrink=10]
% \frametitle{Varianza del estimador MCO: Modelo múltiple}
% Intuición de la varianza del $\hat{\beta}_1$:
% \[
% \text{Var}(\hat{\beta}_1|X) = \frac{\sigma^2}{\text{SST}_1(1 - R_{1}^{2})}
% \]
% \begin{itemize}
%     \item $\text{Var}(\hat{\beta}_1|X)$ es creciente en $\sigma^2$.
%     \begin{itemize}
%         \item Mientras más ruido no explicado haya en la ecuación, más difícil es estimar el efecto parcial de $x_1$ en $y$.
%     \end{itemize}
%     \item $\text{Var}(\hat{\beta}_1|X)$ es decreciente en $\text{SST}_1$.
%     \begin{itemize}
%         \item Esto implica que se prefiere una mayor variación muestral en $x_1$.
%         \item Una alternativa para aumentar esta cantidad y reducir la varianza es aumentar el tamaño muestral.
%         \item El supuesto 3 asegura que $\text{SST}_1 \neq 0$.
%     \end{itemize}
%     \item $\text{Var}(\hat{\beta}_1|X)$ es creciente en $R_{1}^{2}$.
%     \begin{itemize}
%         \item Intuición: Meter más variables le va quitando variabilidad a $x_1$. Esto hace que su estimador sea más impreciso.
%     \end{itemize}
% \end{itemize}
% \end{frame}

% \begin{frame}[shrink=10]
% \frametitle{Problema de Multicolinealidad}
% \[
% \text{Var}(\hat{\beta}_1|X) = \frac{\sigma^2}{SST_1(1 - R^2_1)}
% \]
% \begin{itemize}
%     \item Si dos variables son linealmente dependientes o están muy correlacionadas, la varianza del estimador será muy alta. En el límite, $R^2_1 = 1$ y la varianza es infinita. Esto viola el supuesto 3 de no multicolinealidad perfecta. En este caso, hay multicolinealidad perfecta.
%     \item Trade-offs:
%     \begin{enumerate}
%         \item \textbf{Varianza:} Incluir más variables explicativas \textcolor{red}{reduce la SSR ($\sigma^2$)}, pero \textcolor{blue}{aumenta $R^2_1$}. Efecto ambiguo sobre la varianza $\hat{\text{Var}}(\hat{\beta}_1)$.
%         \item \textbf{Sesgo:} Incluir más variables puede evitar el sesgo del coeficiente si la variable omitida correlaciona con la $x$. A menudo sacrificamos aumentar la varianza para estar seguros de evitar el sesgo.
%     \end{enumerate}
% \end{itemize}
% \end{frame}
\end{document}